\chapter{The Web via Emscripten}
\label{ch:web-emscripten}

CNA's Web target compiles the same C\texttt{++}23 codebase through Emscripten rather than a
native toolchain, and this chapter covers the specific build-system machinery that makes that
possible, along with an honest account of exactly how far it has actually been verified.

\section{Build-system adjustments specific to Emscripten}

Two adjustments apply project-wide before any backend-specific configuration happens.
Emscripten disables C\texttt{++} exception unwinding by default, for code-size reasons; CNA
force-re-enables it globally, since exceptions are used throughout the codebase (every
\texttt{NotImplementedException}, every \texttt{ArgumentOutOfRangeException} this book has
cited depends on real unwinding working). And \texttt{EASYGL} (targeting WebGL2, which is
GLES~3.0) is the default backend under Emscripten, the same default it holds on Linux ---
while \texttt{CANVAS} (Chapter~\ref{ch:canvas-ascii-backends}) is hard-gated the other
direction, refusing to configure at all outside Emscripten, since HTML \texttt{<canvas>} has
no meaning anywhere else.

\begin{lstlisting}[style=cnashell]
# Worked example: the exact compile/link options the top-level CMakeLists.txt
# applies globally under Emscripten, before sharp-runtime is even added --
# unwinding requires every frame in the call chain to be compiled with
# -fexceptions, so this has to happen before anything else is configured.
# (from CMakeLists.txt's own if(EMSCRIPTEN) block)
add_compile_options(-fexceptions)
add_link_options(-fexceptions -sNO_DISABLE_EXCEPTION_CATCHING=1)
\end{lstlisting}

Several test categories are excluded specifically under Emscripten, each for a concrete
reason rather than a blanket ``skip on web'' rule: pixel/integration test suites for every
backend's own CMake test registration are excluded here, and a handful of process-spawning
tests (a two-process networking loopback test, a gamer-services-dispatcher hang-regression
test, an audio mixer test, a glTF-tool test) are excluded specifically because there is no
real multi-process spawning available inside a single Node.js/WebAssembly module.

\section{Making the shared test binary actually exit, and actually yield}

Two further Emscripten-specific link flags apply only to the shared \texttt{CnaTests} binary,
each solving a distinct, concrete problem. Emscripten's default behavior keeps its JavaScript
runtime alive after \texttt{main()} returns, so \texttt{node CnaTests.js} would otherwise never
exit on its own --- a runtime-exit flag corrects this specifically for the test binary.
Separately, the system-link networking loopback tests need to actually yield control back to
Node's own event loop mid-test (so a WebSocket handshake, for instance, has a chance to
complete) --- ordinary synchronous, single-threaded Emscripten execution cannot yield at all
without help, confirmed directly: a real one-second sleep loop never let a WebSocket handshake
finish without Emscripten's Asyncify transformation enabled, so Asyncify is applied, scoped
specifically to this one test binary rather than the whole project.

\begin{lstlisting}[style=cnashell]
# Worked example: both CnaTests-only link flags, exactly as the real
# cmake/UnitTests.cmake applies them -- neither affects any other
# Emscripten target (demos, the two-process harness) in the same build.
target_link_options(CnaTests PRIVATE -sEXIT_RUNTIME=1)
target_link_options(CnaTests PRIVATE -sASYNCIFY=1)
\end{lstlisting}

\section{What has actually been built and run, versus what has not}

This is the section most worth reading carefully, because the picture differs sharply by
backend. The \texttt{CANVAS} backend has a real, confirmed success story: a genuine
\texttt{emcmake} / \texttt{emcc} (version 6.0.2) configure-and-build succeeded, producing a real
build directory for the first time in this project's history, and \texttt{CnaTests.js}
genuinely links and runs under \texttt{node}, with a real, backend-agnostic GoogleTest suite
passing in full.

\begin{lstlisting}[style=cnashell]
# Worked example: the real invocation shape behind that CANVAS success
# story, following the same -DCNA_GRAPHICS_BACKEND convention this book's
# Building CNA chapter already established for native builds.
emcmake cmake -S . -B build-canvas -DCNA_GRAPHICS_BACKEND=CANVAS
cmake --build build-canvas
node build-canvas/tests/CnaTests.js
# The GoogleTest suite genuinely passes here -- but see the next section
# for the hard ceiling on what "passes" can actually prove in this
# environment, even for a build that genuinely links and runs.
\end{lstlisting}

The default \texttt{EASYGL}/WebGL2 path tells a different, more honest story: at the time of
writing, no Emscripten build directory has ever existed for \texttt{EASYGL} at all, and no
automated workflow builds it under Emscripten either. Every specific claim in this project's
own Web/Emscripten graphics-limitations documentation about \texttt{EASYGL} behavior under
WebGL2 --- geometry and tessellation shader absence, anisotropic-filtering extension gating,
whether a WebGL1 fallback exists --- is explicitly labeled a \emph{design-time expectation},
not something empirically verified, precisely because the build that would let someone verify
it has never actually been produced.

\subsection{One real asset the \texttt{CANVAS} build left behind: a reusable SDL3-for-Emscripten cache}

The \texttt{CANVAS} build described above did not just prove \texttt{emcc} works in this
project's dev environment; it also left a concrete, still-present artifact on disk that
directly affects the cost of the next attempt at \texttt{EASYGL} under Emscripten. CNA's
vendored-SDL build machinery (\texttt{cmake/ThirdPartySDL.cmake}) installs SDL3, SDL3\_image,
and SDL3\_mixer once into a persistent cache directory keyed by target platform ---
\texttt{.sdl-prebuilt-emscripten/} specifically for Emscripten, distinct from the
per-native-platform \texttt{.sdl-prebuilt-Linux-x86\_64} / \texttt{.sdl-prebuilt-Windows-x86\_64}
directories this book's Building CNA chapter (Chapter~\ref{ch:building-cna}) already
describes --- and the cache's own comment states its purpose plainly: it lives
\emph{``OUTSIDE any cmake build tree, so \texttt{cmake --build --clean-first} (or deleting the
build directory entirely) does NOT remove SDL artefacts.''} Deleting an Emscripten build
directory and reconfiguring from scratch does not force SDL3 to rebuild; only manually
deleting \texttt{.sdl-prebuilt-emscripten/} itself does that.

That cache is real and still present: a genuine WebAssembly static-library build of all three
SDL3 libraries exists on disk right now, confirmed directly by inspecting the installed
archives rather than trusting a build log --- \texttt{libSDL3.a} at 2{,}024{,}910 bytes,
\texttt{libSDL3\_image.a} at 321{,}570 bytes, \texttt{libSDL3\_mixer.a} at 384{,}298 bytes, plus
a \texttt{libSDL3\_test.a} at 128{,}890 bytes, alongside a proper CMake package-config
install (\texttt{install/lib/cmake/\{SDL3,SDL3\_image,SDL3\_mixer\}/}) that lets any later
configure step find them via ordinary \texttt{find\_package()} rather than hand-wired paths.
Crucially, this build machinery is entirely backend-agnostic: \texttt{ThirdPartySDL.cmake}
branches only on platform (\texttt{EMSCRIPTEN} / \texttt{ANDROID} / \texttt{WIN32} / \texttt{APPLE}/
plain Unix), never on \texttt{CNA\_GRAPHICS\_BACKEND}. That means the SDL3-for-Emscripten cache
the \texttt{CANVAS} build populated is exactly the same cache an \texttt{EASYGL} Emscripten
build would consume --- whoever attempts that build next does not start from zero; the
slowest, most failure-prone part of a first Emscripten configure (compiling SDL3 and its two
satellite libraries to WebAssembly) is already done and sitting on disk, and only the
\texttt{EASYGL}-specific and \texttt{easy-gl}-sibling-repository code paths remain genuinely
unattempted.

\section{Why even a successful build cannot be pixel-verified here}

Even where a build genuinely succeeds and a test binary genuinely runs, there is a hard
ceiling on what it can prove: \texttt{node} provides no WebGL context and no
\texttt{CanvasRenderingContext2D}/DOM at all --- confirmed directly, since
\texttt{SDL\_Init(SDL\_INIT\_VIDEO)} itself throws under Emscripten running inside
\texttt{node}, and the project's own build comments record the exact underlying error rather
than a paraphrase: \texttt{"window is not defined"}, the literal exception a real, genuine
\texttt{CANVAS} smoke-test build hits the moment it tries to touch a browser global that
simply does not exist in a bare Node.js process. This is the same ceiling Chapter~\ref{ch:canvas-ascii-backends} already
described for \texttt{CANVAS} specifically: any test suite that constructs a real
\cnaclass{GraphicsDevice} simply cannot execute in this dev loop at all, real browser or not,
which is exactly why \texttt{docs/canvas-backend.md}'s own manual browser verification
checklist --- covering rotation and origin math, blend modes, and wrap/mirror sampling among
thirteen total items --- is left entirely unchecked, deferred to a human running the game in
an actual browser via a tool like \texttt{emrun}, rather than claimed as done.

One further, genuinely interesting piece of real code is worth knowing about even though it
has never been exercised against a real browser: \texttt{EASYGL}'s own backend implementation
contains real, reviewed WebGL-context-loss handling --- registering for the browser's own
\texttt{WEBGL\_lose\_context} extension and its \texttt{webglcontextlost} /
\texttt{webglcontextrestored} events. This code path exists and has been read and reviewed,
but --- consistent with everything else in this section --- has never actually run against a
real browser; only the ordinary desktop code path (a synchronous SDL GL context
destroy/recreate) has ever actually executed.

\subsection{What the WebGL-context-loss code actually does, read directly}

\texttt{EasyGLGraphicsBackend.cpp} defines two \texttt{EM\_JS} blocks --- real, inline
JavaScript compiled directly into the Emscripten build, not a C\texttt{++} approximation of
browser behavior --- gated entirely behind \texttt{\#if defined(\_\_EMSCRIPTEN\_\_)}, so none
of this code exists at all in a native desktop build:

\begin{lstlisting}
EM_JS(void, CNA_DebugLoseWebGLContext, (), {
    const canvas = Module['canvas'] || document.querySelector('canvas');
    const gl = Module['ctx'] || canvas.getContext('webgl2') || canvas.getContext('webgl');
    const ext = gl.getExtension('WEBGL_lose_context');
    ext.loseContext();
});
// CNA_DebugRestoreWebGLContext mirrors this exactly, calling ext.restoreContext() instead.
\end{lstlisting}

Both functions look up the live \texttt{WebGLRenderingContext} the same way (preferring
Emscripten's own \texttt{Module['ctx']} handle, falling back to querying the canvas directly
for a \texttt{webgl2} or \texttt{webgl} context), then call the real
\texttt{WEBGL\_lose\_context} extension's \texttt{loseContext()} / \texttt{restoreContext()}
methods --- the same browser mechanism a real GPU driver crash or a backgrounded tab actually
triggers, not a simulation built from first principles. \texttt{metagl::\allowbreak{}Install\allowbreak{}Emscripten\allowbreak{}ContextLossCallbacks()},
installed once at backend construction, is what turns the resulting
\texttt{webglcontextlost} / \texttt{webglcontextrestored} canvas events into calls to
\texttt{metagl::NotifyContextLost()} / \texttt{NotifyContextRestored()} --- the same two
notification points the desktop code path calls directly and synchronously.

That synchronicity difference is the one genuine architectural divergence between the two
platforms, and it is handled explicitly rather than papered over. \texttt{DebugSimulateContextLoss()}'s
own source comment states it plainly for the Emscripten branch: ``The \texttt{webglcontextlost}
canvas event fires asynchronously and triggers \texttt{metagl::NotifyContextLost()} via
\texttt{InstallEmscriptenContextLossCallbacks()}'' --- the call to
\texttt{CNA\_DebugLoseWebGLContext()} only \emph{requests} the loss; the actual
\cnaclass{ResourceRegistry} notification happens later, asynchronously, whenever the browser's
own event loop delivers the event. The desktop \texttt{\#else} branch is the opposite shape
entirely: a four-step synchronous sequence performed inline in the same function call ---
\texttt{metagl::NotifyContextLost()} (releasing every tracked GL handle without issuing GL
calls, since the context is still technically valid at that point), resetting every stock-effect
shader program handle via \texttt{reset\_no\_gl()}, destroying and recreating the real SDL GL
context (\texttt{SDL\_GL\_DestroyContext} / \texttt{SDL\_GL\_CreateContext}), reloading every GL
function pointer, and finally \texttt{metagl::NotifyContextRestored()} to make every tracked
resource recreate itself. \texttt{DebugRestoreContext()} reflects the same asymmetry one level
up: on Emscripten it is a second, independent call (mirroring the browser's own two separate
events); on desktop it is simply \texttt{DebugSimulateContextLoss()} called a second time,
since loss and restore are one atomic operation there with no separate event to wait for.

\subsection{A real, still-open inconsistency: only one of three Emscripten targets pins a WebGL version}

Reading \texttt{cmake/Examples.cmake} directly turns up a genuine, previously-undocumented gap
in this project's own Emscripten link-flag consistency. \texttt{cna\_house3d\_demo} --- the one
real 3D EasyGL demo --- explicitly pins its WebGL version under Emscripten:

\begin{lstlisting}[style=cnashell]
target_link_options(cna_house3d_demo PRIVATE
    -sALLOW_MEMORY_GROWTH=1
    -sSTACK_SIZE=1048576
    -sINITIAL_MEMORY=67108864
    -sFORCE_FILESYSTEM=1
    "-sMIN_WEBGL_VERSION=2"         # WebGL 2 = OpenGL ES 3.0
    "-sMAX_WEBGL_VERSION=2"
)
\end{lstlisting}

\texttt{cna\_demo\_2d}'s own Emscripten \texttt{target\_link\_options} block, confirmed by
reading it directly rather than assuming symmetry with the 3D demo, has no equivalent pin at
all: \texttt{-sALLOW\_MEMORY\_GROWTH=1}, a smaller \texttt{-sSTACK\_SIZE}, a larger
\texttt{-sINITIAL\_MEMORY} (134~MB versus 64~MB, sized for its preloaded \texttt{Content}
directory rather than a procedural scene), \texttt{-sFORCE\_FILESYSTEM=1}, and a
\texttt{--preload-file} flag for that content --- but no \texttt{MIN\_WEBGL\_VERSION} /
\texttt{MAX\_WEBGL\_VERSION} anywhere in the block. Without the pin, Emscripten is free to
negotiate WebGL~1 with a browser/GPU combination that does not offer WebGL~2, silently handing
that build a GLES~2-class context instead of the GLES~3.0-class one \texttt{EasyGL}'s own code
assumes everywhere else. This is exactly the failure mode \texttt{docs/web-emscripten-graphics-limitations.md}
names as a documented risk rather than a hypothetical one, and this book's own source read
confirms the doc's claim precisely: \texttt{cna\_demo\_sound} (audio-only, no 3D rendering
surface at all) has the same gap for the same underlying reason --- it was never built with a
3D pipeline in mind, so nobody had reason to add the pin --- but \texttt{cna\_demo\_2d}
genuinely does draw through \cnaclass{SpriteBatch}, meaning it is the one target where an
un-pinned WebGL~1 negotiation could actually produce a silently degraded, unverified rendering
path rather than simply being irrelevant.

\section{The honest summary}

Web support in CNA is real in the sense that a specific, narrow slice of it (\texttt{CANVAS},
building and linking under Node) has been genuinely produced and run; it is aspirational in
the sense that the backend most users would actually reach for (\texttt{EASYGL}/WebGL2) has
never been built under Emscripten at all in this project's own history, and no claim about
its WebGL2-specific behavior should be read as more than a design-time expectation until
someone actually produces that build and checks it in a real browser.
